% questions.tex —— 2011 年·湖北卷·文科 · 题目内容片段（无答案）
% 由 tools/generate_projects.py 依来源总账自 2011/2011_数学真题_湖北文_试题.md 生成，勿手改（改源 MD 后重跑）

\gkpart{一、选择题}

\begin{question}
\gkti{1}已知 $U=\{1,2,3,4,5,6,7,8\}$，$A=\{1,3,5,7\}$，$B=\{2,4,5\}$，则 $\complement_U(A\cup B)=$\mbox{（　　）}

（A）$\{6,8\}$

（B）$\{5,7\}$

（C）$\{4,6,7\}$

（D）$\{1,3,5,6,8\}$

\end{question}
\begin{question}
\gkti{2}若向量 $\boldsymbol a=(1,2)$，$\boldsymbol b=(1,-1)$，则 $2\boldsymbol a+\boldsymbol b$ 与 $\boldsymbol a-\boldsymbol b$ 的夹角等于\mbox{（　　）}

（A）$-\frac\pi4$

（B）$\frac\pi6$

（C）$\frac\pi4$

（D）$\frac{3\pi}4$

\end{question}
\begin{question}
\gkti{3}若定义在 $\mathbf R$ 上的偶函数 $f(x)$ 和奇函数 $g(x)$ 满足 $f(x)+g(x)=\mathrm e^x$，则 $g(x)=$\mbox{（　　）}

（A）$\mathrm e^x-\mathrm e^{-x}$

（B）$\frac12(\mathrm e^x+\mathrm e^{-x})$

（C）$\frac12(\mathrm e^{-x}-\mathrm e^x)$

（D）$\frac12(\mathrm e^x-\mathrm e^{-x})$

\end{question}
\begin{question}
\gkti{4}将两个顶点在抛物线 $y^2=2px$（$p>0$）上，另一个顶点是此抛物线焦点的正三角形个数记为 $n$，则\mbox{（　　）}

（A）$n=0$

（B）$n=1$

（C）$n=2$

（D）$n\geqslant3$

\Needspace{0.4\textheight}
\end{question}
\begin{question}
\gkti{5}有一个容量为 $200$ 的样本，其频率分布直方图如图所示，根据样本的频率分布直方图估计，样本数据落在区间 $[10,12)$ 内的频数为\mbox{（　　）}

（A）$18$

（B）$36$

（C）$54$

（D）$72$
\nopagebreak[4]
\begin{center}% 2011 年湖北卷文科第 5 题：容量 200 的样本频率分布直方图。
% 来源 PDF：普通高考/2011/2011湖北文.pdf，SHA-256 97bda477ec8baea4f3f31a9b6013a77e9a4bcd14f462013754d52c02ff5d1770。
% 原页：第 1 页第 1 栏中部，书页码 669。各组距为 2。
\begin{tikzpicture}[>=stealth, semithick, font=\small, x=0.62cm, y=3.2cm]
  \draw[->] (0,0) -- (13,0) node[right] {样本数据};
  \draw[->] (0,0) -- (0,0.22) node[above] {$\dfrac{\text{频率}}{\text{组距}}$};
  \foreach \x/\h in {1/0.02,3/0.02,5/0.05,7/0.15,9/0.19,11/0.09}{
    \draw (\x-1,0) rectangle (\x+1,\h);
  }
  \foreach \x in {0,2,4,6,8,10,12} \node[below] at (\x,0) {\x};
  \foreach \y in {0.02,0.05,0.15,0.19} \node[left] at (0,\y) {\y};
  \node[below left] at (0,0) {$O$};
\end{tikzpicture}
\end{center}
\nopagebreak[4]
\end{question}
\begin{question}
\gkti{6}已知函数 $f(x)=\sqrt3\sin x-\cos x$，$x\in\mathbf R$。若 $f(x)\geqslant1$，则 $x$ 的取值范围为\mbox{（　　）}

（A）$\left\{x\ \middle|\ 2k\pi+\frac\pi3\leqslant x\leqslant2k\pi+\pi,\ k\in\mathbf Z\right\}$

（B）$\left\{x\ \middle|\ k\pi+\frac\pi3\leqslant x\leqslant k\pi+\pi,\ k\in\mathbf Z\right\}$

（C）$\left\{x\ \middle|\ 2k\pi+\frac\pi6\leqslant x\leqslant2k\pi+\frac{5\pi}6,\ k\in\mathbf Z\right\}$

（D）$\left\{x\ \middle|\ k\pi+\frac\pi6\leqslant x\leqslant k\pi+\frac{5\pi}6,\ k\in\mathbf Z\right\}$

\end{question}
\begin{question}
\gkti{7}设球的体积为 $V_1$，它的内接正方体的体积为 $V_2$，下列说法中最合适的是\mbox{（　　）}

（A）$V_1$ 比 $V_2$ 大约多一半

（B）$V_1$ 比 $V_2$ 大约多两倍半

（C）$V_1$ 比 $V_2$ 大约多一倍

（D）$V_1$ 比 $V_2$ 大约多一倍半

\end{question}
\begin{question}
\gkti{8}直线 $2x+y-10=0$ 与不等式组 $\begin{cases}x\geqslant0,\\y\geqslant0,\\x-y\geqslant-2,\\4x+3y\leqslant20\end{cases}$ 表示的平面区域的公共点有\mbox{（　　）}

（A）$0$ 个

（B）$1$ 个

（C）$2$ 个

（D）无数个

\end{question}
\begin{question}
\gkti{9}《九章算术》“竹九节”问题：现有一根 $9$ 节的竹子，自上而下各节的容积成等差数列，上面 $4$ 节的容积共 $3$ 升，下面 $3$ 节的容积共 $4$ 升，则第 $5$ 节的容积为\mbox{（　　）}

（A）$1$ 升

（B）$\frac{67}{66}$ 升

（C）$\frac{47}{44}$ 升

（D）$\frac{37}{33}$ 升

\end{question}
\begin{question}
\gkti{10}若实数 $a$，$b$ 满足 $a\geqslant0$，$b\geqslant0$，且 $ab=0$，则称 $a$ 与 $b$ 互补。记 $\varphi(a,b)=\sqrt{a^2+b^2}-a-b$，那么 $\varphi(a,b)=0$ 是 $a$ 与 $b$ 互补的\mbox{（　　）}

（A）必要而不充分条件

（B）充分而不必要条件

（C）充要条件

（D）既不充分也不必要条件

\end{question}
\gkpart{二、填空题}

\begin{question}
\gkti{11}某市有大型超市 $200$ 家、中型超市 $400$ 家、小型超市 $1400$ 家。为掌握各类超市的营业情况，现按分层抽样方法抽取一个容量为 $100$ 的样本，应抽取中型超市 $\underline{\qquad}$ 家。

\end{question}
\begin{question}
\gkti{12}$\left(x-\dfrac1{3\sqrt x}\right)^{18}$ 的展开式中含 $x^{15}$ 的项的系数为 $\underline{\qquad}$。（结果用数值表示）

\end{question}
\begin{question}
\gkti{13}在 $30$ 瓶饮料中，有 $3$ 瓶已过了保质期。从这 $30$ 瓶饮料中任取 $2$ 瓶，则至少取到 $1$ 瓶已过保质期饮料的概率为 $\underline{\qquad}$。（结果用最简分数表示）

\end{question}
\begin{question}
\gkti{14}过点 $(-1,-2)$ 的直线 $l$ 被圆 $x^2+y^2-2x-2y+1=0$ 截得的弦长为 $\sqrt2$，则直线 $l$ 的斜率为 $\underline{\qquad}$。

\end{question}
\begin{question}
\gkti{15}里氏震级 $M$ 的计算公式为：$M=\lg A-\lg A_0$，其中 $A$ 是测震仪记录的地震曲线的最大振幅，$A_0$ 是相应的标准地震的振幅。假设在一次地震中，测震仪记录的最大振幅是 $1000$，此时标准地震的振幅为 $0.001$，则此次地震的震级为 $\underline{\qquad}$ 级；$9$ 级地震的最大的振幅是 $5$ 级地震最大振幅的 $\underline{\qquad}$ 倍。

\end{question}
\gkpart{三、解答题}

\begin{question}
\gkti{16}设 $\triangle ABC$ 的内角 $A$，$B$，$C$ 所对的边分别为 $a$，$b$，$c$。已知 $a=1$，$b=2$，$\cos C=\dfrac14$。

（1） 求 $\triangle ABC$ 的周长；

（2） 求 $\cos(A-C)$ 的值。

\end{question}
\begin{question}
\gkti{17}成等差数列的三个正数的和等于 $15$，并且这三个数分别加上 $2$，$5$，$13$ 后成为等比数列 $\{b_n\}$ 中的 $b_3$，$b_4$，$b_5$。

（1） 求数列 $\{b_n\}$ 的通项公式；

（2） 数列 $\{b_n\}$ 的前 $n$ 项和为 $S_n$，求证：数列 $\left\{S_n+\dfrac54\right\}$ 是等比数列。

\Needspace{0.4\textheight}
\end{question}
\begin{question}
\gkti{18}如图，已知正三棱柱 $ABC-A_1B_1C_1$ 的底面边长为 $2$，侧棱长为 $3\sqrt2$，点 $E$ 在侧棱 $AA_1$ 上，点 $F$ 在侧棱 $BB_1$ 上，且 $AE=2\sqrt2$，$BF=\sqrt2$。

（1） 求证：$CF\perp C_1E$；

（2） 求二面角 $E-CF-C_1$ 的大小。
\nopagebreak[4]
\begin{center}% 2011 年湖北卷文科第 18 题：正三棱柱 ABC-A1B1C1，E 在 AA1、F 在 BB1。
% 来源 PDF：普通高考/2011/2011湖北文.pdf，SHA-256 97bda477ec8baea4f3f31a9b6013a77e9a4bcd14f462013754d52c02ff5d1770。
% 原页：第 1 页第 3 栏中部，书页码 669。底面边长 2，侧棱 3√2。
\begin{tikzpicture}[>=stealth, semithick, font=\small, scale=0.95]
  \coordinate (A) at (0,0);
  \coordinate (B) at (3.0,0);
  \coordinate (C) at (1.5,1.0);
  \coordinate (A1) at (0,3.0);
  \coordinate (B1) at (3.0,3.0);
  \coordinate (C1) at (1.5,4.0);
  \coordinate (E) at ($(A)!0.667!(A1)$);
  \coordinate (F) at ($(B)!0.333!(B1)$);
  \draw[thick] (A)--(B)--(C)--cycle;
  \draw[thick] (A1)--(B1)--(C1)--cycle;
  \draw[thick] (A)--(A1) (B)--(B1) (C)--(C1);
  \draw (C)--(F);
  \draw (C1)--(E);
  \draw (E)--(F);
  \fill (E) circle (1.2pt) node[left] {$E$};
  \fill (F) circle (1.2pt) node[right] {$F$};
  \node[left] at (A) {$A$};
  \node[right] at (B) {$B$};
  \node[below] at (C) {$C$};
  \node[left] at (A1) {$A_1$};
  \node[right] at (B1) {$B_1$};
  \node[above] at (C1) {$C_1$};
\end{tikzpicture}
\end{center}
\nopagebreak[4]
\end{question}
\begin{question}
\gkti{19}提高过江大桥的车辆通行能力可改善整个城市的交通状况。在一般情况下，大桥上的车流速度 $v$（单位：千米/小时）是车流密度 $x$（单位：辆/千米）的函数。当桥上的车流密度达到 $200$ 辆/千米时，造成堵塞，此时车流速度为 $0$；当车流密度不超过 $20$ 辆/千米时，车流速度为 $60$ 千米/小时。研究表明：当 $20\leqslant x\leqslant200$ 时，车流速度 $v$ 是车流密度 $x$ 的一次函数。

（1） 当 $0\leqslant x\leqslant200$ 时，求函数 $v(x)$ 的表达式；

（2） 当车流密度 $x$ 为多大时，车流量（单位时间内通过桥上某观测点的车辆数，单位：辆/小时）$f(x)=x\cdot v(x)$ 可以达到最大，并求出最大值。（精确到 $1$ 辆/小时）

\end{question}
\begin{question}
\gkti{20}设函数 $f(x)=x^3+2ax^2+bx+a$，$g(x)=x^2-3x+2$，其中 $x\in\mathbf R$，$a$，$b$ 为常数，已知曲线 $y=f(x)$ 与 $y=g(x)$ 在点 $(2,0)$ 处有相同的切线 $l$。

（1） 求 $a$，$b$ 的值，并写出切线 $l$ 的方程；

（2） 若方程 $f(x)+g(x)=mx$ 有三个互不相同的实根 $0$，$x_1$，$x_2$，其中 $x_1<x_2$，且对任意的 $x\in[x_1,x_2]$，$f(x)+g(x)<m(x-1)$ 恒成立，求实数 $m$ 的取值范围。

\end{question}
\begin{question}
\gkti{21}平面内与两定点 $A_1(-a,0)$，$A_2(a,0)$（$a>0$）连线的斜率之积等于非零常数 $m$ 的点的轨迹，加上 $A_1$，$A_2$ 两点所成的曲线 $C$ 可以是圆、椭圆或双曲线。

（1） 求曲线 $C$ 的方程，并讨论 $C$ 的形状与 $m$ 值的关系；

（2） 当 $m=-1$ 时，对应的曲线为 $C_1$；对给定的 $m\in(-1,0)\cup(0,+\infty)$，对应的曲线为 $C_2$。设 $F_1$，$F_2$ 是 $C_2$ 的两个焦点。试问：在 $C_1$ 上，是否存在点 $N$，使得 $\triangle F_1NF_2$ 的面积 $S=|m|a^2$？若存在，求 $\tan\angle F_1NF_2$ 的值；若不存在，请说明理由。

\end{question}